\documentclass[a4paper]{article}

\usepackage[utf8x]{inputenc}    
\usepackage[T1]{fontenc}
\usepackage{graphicx}
\graphicspath{{Figures/}}
\usepackage[frenchb]{babel}
\usepackage[francais,bloc,nowatermark]{automultiplechoice}    
% [nowatermark] pour enlever le texte "PROJET" et le pied de page
\usepackage{multicol}
\usepackage{tkz-tab}
\usepackage{fancyhdr,amssymb,amsmath}
%\DeclareGraphicsRule{.7}{mps}{*}{}

%===================
\begin{document}
%===================

%%% On répartie les quesitons dans différents groupes.
% ---------------------------------
%%% Questions groupe : "catégorie1"
%----------------------------------
\element{categorie1}{
  \begin{question}{Derive1}    
% Question ouverte   
\begin{enumerate}
\item Développer, réduire et ordonner :
\begin{enumerate}
\begin{multicols}{2}
\item $f(x)=(2x-3)(x+1)$
\item $g(x)=-5(x-4)^{2}$
\end{multicols}
\end{enumerate}
\item Montrer que $-1$ est une racine de $f$.
\end{enumerate}

   \AMCOpen{lines=5,dots=false}{\wrongchoice{NA}\scoring{0}\wrongchoice{AB}\scoring{1}\wrongchoice{B}\scoring{2}\correctchoice{TB}\scoring{3}}
% AB rapporte 1 point
% B rapporte 2 points
% TB rapporte 3 points     
  \end{question}
}

%
\element{categorie2}{
  \begin{question}{Derive2}  
% Question ouverte  
\begin{enumerate}
\item Résoudre dans $\mathbb{R}$ les équations suivantes. On donnera la valeur exacte puis une valeur approchée à 0,01 près des solutions.

\begin{enumerate}
\begin{multicols}{4}
\item $x^{2}=100$
\item $2 x^{2}+1=5$
\item $x^{3}=1000$
\item $2 x^{3}+11=5$
\end{multicols}
\end{enumerate}

\item Soit $f$ la fonction définie sur $\mathbb{R}$ par :
$f(x)=-3(x-8)(x+5)(x-3)$
\begin{enumerate}
	\item Justifier que $f$ est un polynôme de degré 3.
	\item Déterminer les racines de $f$ \textit{(à l'aide de sa forme factorisée)}.
	\item Dresser le tableau de signes de $f$ sur $\mathbb{R}$.
	\item En déduire les solutions de l'inéquation $f(x)<0$.	
\end{enumerate}
\end{enumerate}
% AB rapporte 1 point
% B rapporte 2 points
% TB rapporte 3 points         
  \end{question}
   \AMCOpen{lines=16,dots=false}{\wrongchoice{NA}\scoring{0}\wrongchoice{AB}\scoring{1}\wrongchoice{B}\scoring{2}\correctchoice{TB}\scoring{3}}
}




\element{categorie4}{
  \begin{question}{variation1}  
% Question ouverte  
On a représenté graphiquement ci-dessous une fonction $f$ polynôme du troisième degré sur un intervalle et telle que $f(x)=a(x-x_1)(x-x_2)(x-x_3)$
\begin{enumerate}
\item Déterminer graphiquement le tableau de variation de $f$ sur l'intervalle $[-1;2]$.
\item Quelles sont les valeurs de $x_{1}, x_{2}$ et $x_{3}$ ?
\item Déterminer la valeur de $a$ sachant que $f(0)=1$.
\item En déduire l'expression de $f(x)$.
\end{enumerate}
\vspace*{-2cm}
\hspace*{7cm}\includegraphics[scale=1]{courbeax2b.9}
\vspace{2cm}

   \AMCOpen{lines=10,dots=false}{\wrongchoice{NA}\scoring{0}\wrongchoice{AB}\scoring{1}\wrongchoice{B}\scoring{2}\correctchoice{TB}\scoring{3}}
% AB rapporte 1 point
% B rapporte 2 points
% TB rapporte 3 points         
  \end{question}
}


\element{categorie4}{
  \begin{question}{variation2}  
% Question ouverte  
On considère la fonction $f$ définie sur l'intervalle $[-4 ; 4]$ par $f(x)=x^{2}-2 x-3$.
\begin{enumerate}
\item  Calculer l'image de -1 par $f$.
\item  Montrer que 3 est solution de l'équation $f(x)=0$.
\item  En utilisant les questions 1 et 2, donner une forme factorisée de $f(x)$.
\item Parmi les trois courbes suivantes, déterminer, en justifiant, celle qui représente graphiquement la fonction $f$.
\end{enumerate}
\vspace*{-3.5cm}
\hspace*{6cm}\includegraphics[scale=0.7]{courbeax2b.12}
\vspace{3.2cm}

   \AMCOpen{lines=8,dots=false}{\wrongchoice{NA}\scoring{0}\wrongchoice{AB}\scoring{1}\wrongchoice{B}\scoring{2}\correctchoice{TB}\scoring{3}}
% AB rapporte 1 point
% B rapporte 2 points
% TB rapporte 3 points         
  \end{question}
}

%-------------------------

%%%%%%%%%%%%%%%%%%%%%%%%%
%% fin des questions %%%%
%%%%%%%%%%%%%%%%%%%%%%%%%


%%%%%%%%%%%%%%%%%%%%%%%%%%%%
%%% fabrication des copies%%
%%%%%%%%%%%%%%%%%%%%%%%%%%%%
\exemplaire{1}{    

%%% debut de l'entête des copies :    
	%%% première ligne

	\begin{minipage}{1\linewidth}
	  \centering\large\bf DS de Mathématiques 1STMG sur les polynômes   
	\end{minipage}

	%%% codage des numéros d'étudiants sur 3 chiffres par les candidats et encadré d'écriture manuelle des noms.
%	\vspace*{3mm}
	{
	%\setlength{\parindent}{0pt}\AMCcode{etu}{2}\hspace*{\fill}
	\begin{minipage}[b]{12cm}
	% \hspace*{5mm}
  	%  Codez votre numéro ci contre chiffre par chiffre, puis complétez l'encadré.
    % \vspace{3ex}

	 \hfill\champnom{
	 \centering
	 \fbox{
	 \begin{minipage}{1\linewidth}
		\vspace*{3mm}
		NOM - Prénom - Classe :
		\vspace*{3mm}
	 \end{minipage}
	    }	
	   }
	\hfill\vspace{1ex}
% texte de présentation 
\begin{center}\em

Durée : 55 minutes.

  Aucun document n'est autorisé. La calculatrice est autorisée.\\
  Toutes les réponses doivent être justifiées.
%  Les questions faisant apparaître le symbole \multiSymbole{} peuvent
%  présenter une ou plusieurs bonnes réponses. Les autres ont une unique bonne réponse.
%  Les réponses fausses ou incohérentes retirent des points. Pour éviter les erreurs machines, il est nécessaire de \textbf{noircir vos cases}.

\end{center}

	\end{minipage}\hspace*{\fill}
	}

\vspace{4ex}

%%% Fin de l'entête
%%%%%%%%%%%%%%%%%%%%%%%

%% Composition des QCM
%%% On mélange des questions par groupe
% On efface le {tout} précédent
\cleargroup{tout}

% On mélange par "catégorie", et on prélève aléatoirement [n] questions 
% du lot {categorie-k} vers la copie complète {tout}
% puis on restitue la copie.
% Si on met [0], toutes les questions sont prises.
\melangegroupe{categorie1}\copygroup[0]{categorie1}{tout}
\melangegroupe{categorie2}\copygroup[0]{categorie2}{tout}
%\melangegroupe{categorie3}\copygroup[0]{categorie3}{tout}
\melangegroupe{categorie4}\copygroup[0]{categorie4}{tout}

\restituegroupe{tout}

% Saut d'une page si le sujet comporte un nombre impair de pages.
% Ainsi il est possible d'imprimer les sujets en recto-verso.

\AMCcleardoublepage    

%%% Fin de la création des questionnaires

}
% Fin de exemplaires{k}{ 

\end{document}

